% !TeX encoding = UTF-8
% !TeX spellcheck = en_US
% !TeX TXS-program:compile = txs:///pdflatex/[--shell-escape]
%
% The Diabolical AI Art Battle: Why LLMs Suck at ASCII Art (v2)
% Franklin Diaz — November 18, 2026
% Transamerica Cybersecurity Club — CU Denver Business School

\documentclass[aspectratio=169,11pt,xcolor={table,dvipsnames}]{beamer}

%% ──────────────────────────────────────────────
%% preamble.tex — High-Contrast Modern Cyber Theme
%% Aligned with v2 Quest for Greebling design language
%% ──────────────────────────────────────────────

\usepackage{graphicx}
\usepackage{amsmath}
\usepackage{qrcode}
\usepackage{xurl}
\usepackage{hyperref}
\usepackage{booktabs}
\usepackage{tikz}
\usetikzlibrary{arrows.meta, positioning}
\usepackage{listings}
\usepackage{xcolor}
\usepackage{helvet}
\renewcommand{\familydefault}{\sfdefault}
\usepackage[T1]{fontenc}
\usepackage[absolute,overlay]{textpos}

% --- HIGH CONTRAST MODERN CYBER THEME PALETTE ---
\definecolor{bgdark}{RGB}{15, 23, 42}           % Deep navy background (#0F172A)
\definecolor{textlight}{RGB}{242, 245, 249}      % Crisp white/light-gray text (#F8FAFC)
\definecolor{textmuted}{RGB}{148, 163, 184}      % Muted secondary text (#94A3B8)
\definecolor{accentcyan}{RGB}{0, 210, 255}       % Electric cyan (#00D2FF)
\definecolor{accentmagenta}{RGB}{217, 70, 239}   % Neon magenta (#D946EF)
\definecolor{accentemerald}{RGB}{0, 230, 118}    % Emerald green (#00E676)
\definecolor{accentamber}{RGB}{255, 179, 0}      % Alert amber (#FFB300)
\definecolor{accentcoral}{RGB}{250, 88, 45}      % Warm orange-red bullet (#FA582D)
\definecolor{deepslate}{RGB}{30, 41, 59}         % Card/box background (#1E293B)
\definecolor{headerblue}{RGB}{12, 17, 26}        % Deep bar header tone
\definecolor{borderblue}{RGB}{51, 65, 85}        % Card borders (#334155)
\definecolor{codebg}{RGB}{10, 15, 28}            % Code & terminal background

% Legacy color aliases for existing slides
\definecolor{bulletcolor}{RGB}{250, 88, 45}
\definecolor{accent1}{RGB}{217, 70, 239}
\definecolor{accent2}{RGB}{139, 92, 246}
\definecolor{accent3}{RGB}{0, 210, 255}
\definecolor{headlight}{RGB}{255, 179, 0}

% --- HYPERLINK SETUP ---
\hypersetup{colorlinks=true,urlcolor=accentcyan,linkcolor=accentmagenta}

% --- BEAMER THEME COLOR OVERRIDES ---
\setbeamercolor{background canvas}{bg=bgdark}
\setbeamercolor{normal text}{fg=textlight}
\setbeamercolor{frametitle}{bg=headerblue, fg=accentcyan}
\setbeamercolor{frametitle right}{bg=headerblue}
\setbeamercolor{title}{bg=headerblue, fg=white}
\setbeamercolor{subtitle}{fg=accentcyan}
\setbeamercolor{author}{fg=textlight}
\setbeamercolor{institute}{fg=textmuted}
\setbeamercolor{date}{fg=accentamber}
\setbeamercolor{item}{fg=accentcyan}
\setbeamercolor{subitem}{fg=accentmagenta}
\setbeamercolor{subsubitem}{fg=accentamber}
\setbeamercolor{description item}{fg=accentcyan}
\setbeamercolor{section in toc}{fg=white}
\setbeamercolor{section in toc shaded}{fg=textmuted}
\setbeamertemplate{itemize items}[circle]

% Footer & Header Colors
\setbeamercolor{author in head/foot}{bg=headerblue, fg=textmuted}
\setbeamercolor{title in head/foot}{bg=deepslate, fg=accentcyan}
\setbeamercolor{date in head/foot}{bg=headerblue, fg=textmuted}

% Block Colors
\setbeamercolor{block title}{bg=deepslate!90!accentcyan, fg=accentcyan}
\setbeamercolor{block body}{bg=deepslate, fg=textlight}
\setbeamercolor{block title alerted}{bg=deepslate!90!accentamber, fg=accentamber}
\setbeamercolor{block body alerted}{bg=deepslate, fg=textlight}
\setbeamercolor{block title example}{bg=deepslate!90!accentmagenta, fg=accentmagenta}
\setbeamercolor{block body example}{bg=deepslate, fg=textlight}

% Navigation & Footline
\setbeamertemplate{navigation symbols}{}
\setbeamertemplate{footline}{
  \leavevmode%
  \hbox{%
    \begin{beamercolorbox}[wd=.333333\paperwidth,ht=2.5ex,dp=1.1ex,left,leftskip=2ex]{author in head/foot}%
      \usebeamerfont{author in head/foot}\textbf{Franklin Diaz} (\texttt{franklin.diaz@ucdenver.edu})
    \end{beamercolorbox}%
    \begin{beamercolorbox}[wd=.333333\paperwidth,ht=2.5ex,dp=1.1ex,center]{title in head/foot}%
      \usebeamerfont{title in head/foot}The Diabolical AI Art Battle
    \end{beamercolorbox}%
    \begin{beamercolorbox}[wd=.333333\paperwidth,ht=2.5ex,dp=1.1ex,right,rightskip=2ex]{date in head/foot}%
      \usebeamerfont{date in head/foot}Nov 18, 2026 \quad $\cdot$ \quad \insertframenumber{} / \inserttotalframenumber%
    \end{beamercolorbox}%
  }%
  \vskip0pt%
}

% Presenter Note Configuration (Clean readable dark-on-light in note view)
\setbeamercolor{note page}{bg=white, fg=black}
\setbeamercolor{note title}{bg=gray!20, fg=black}
\setbeamertemplate{note page}{%
  \begin{minipage}{\textwidth}
    \color{black}
    \insertnote{}
  \end{minipage}
}

% --- CODE LISTINGS STYLE ---
\lstset{
  backgroundcolor=\color{codebg},
  basicstyle=\ttfamily\scriptsize\color{textlight},
  breaklines=true,
  captionpos=b,
  commentstyle=\color{textmuted}\itshape,
  keywordstyle=\color{accentcyan}\bfseries,
  stringstyle=\color{accentemerald},
  numbers=none,
  frame=single,
  rulecolor=\color{borderblue},
  showspaces=false,
  showstringspaces=false,
  tabsize=2
}

% --- SECTION TRANSITION SLIDES ---
\AtBeginSection[]{
  \begin{frame}[plain]
    \vfill
    \centering
    \begin{beamercolorbox}[sep=14pt,center,rounded=true,shadow=false]{frametitle}
      {\Huge \textcolor{accentcyan}{\bfseries \insertsectionhead}}\par%
      \vspace{0.4cm}
      {\small \textcolor{textmuted}{Transamerica Cybersecurity Club \quad$\cdot$\quad CU Denver Cyber Defense Center}}
    \end{beamercolorbox}
    \vfill
  \end{frame}
}

% --- REUSABLE HIGHLIGHT HELPERS ---
\newcommand{\hl}[1]{\textcolor{accentcyan}{\textbf{#1}}}
\newcommand{\hlm}[1]{\textcolor{accentmagenta}{\textbf{#1}}}
\newcommand{\hla}[1]{\textcolor{accentamber}{\textbf{#1}}}
\newcommand{\hle}[1]{\textcolor{accentemerald}{\textbf{#1}}}
\newcommand{\hlg}[1]{\textcolor{accentamber}{\textbf{#1}}}
\newcommand{\hll}[1]{\textcolor{accentcyan}{\textbf{#1}}}


\title[Why LLMs Suck at ASCII Art v2]{The Diabolical AI Art Battle}
\subtitle{Tokenization, Spatial Blind Spots, and the AI Myth}
\author[Franklin Diaz]{Franklin Diaz\texorpdfstring{\\ \small\texttt{franklin.diaz@ucdenver.edu} \quad \texttt{@n0ctilucent}}{}}
\institute[CU Denver]{
  \texorpdfstring{Transamerica Cybersecurity Club \\ Cyber Defense Center (NCAE-C) \\ University of Colorado Denver --- Business School}{Transamerica Cybersecurity Club, CU Denver Business School}
}
\date[Nov 18, 2026]{\texorpdfstring{November 18, 2026 --- 5:00 PM to 6:30 PM MST \\ \small JP Morgan Room (1200), Business School}{November 18, 2026}}

\begin{document}

%%%%%%%%%%%%%%%%%%%%%%%%%%%%%%%%%%%%%%%%%%%%%%%%%%%%%%%%%%%%%%%%
%% TITLE SLIDE
%%%%%%%%%%%%%%%%%%%%%%%%%%%%%%%%%%%%%%%%%%%%%%%%%%%%%%%%%%%%%%%%
\begin{frame}[plain]
  \vfill
  \begin{columns}[c]
    \begin{column}{0.73\textwidth}
      {\Huge \textcolor{white}{\textbf{The Diabolical AI Art Battle}}}\\[0.25cm]
      {\Large \textcolor{accentcyan}{\textbf{Tokenization, Spatial Blind Spots, and the AI Myth}}}\\[0.5cm]
      {\normalsize \textbf{Franklin Diaz} \quad \textcolor{accentamber}{\texttt{franklin.diaz@ucdenver.edu}} \quad \textcolor{textmuted}{\texttt{@n0ctilucent}}}\\[0.18cm]
      {\small \textcolor{textmuted}{Transamerica Cybersecurity Club $\cdot$ Cyber Defense Center (NCAE-C)}}\\[0.08cm]
      {\small \textcolor{textmuted}{University of Colorado Denver --- Business School}}\\[0.35cm]
      {\footnotesize \textcolor{accentamber}{\textbf{Wednesday, November 18, 2026}} $\cdot$ 5:00~PM -- 6:30~PM MST $\cdot$ JP Morgan Room (1200)} % chktex 8
    \end{column}
    \begin{column}{0.24\textwidth}
      \centering
      \colorbox{white}{\qrcode[height=2.5cm]{https://github.com/n0ctilucent/talk-ccdc-nov-26}}\\[0.2cm]
      {\tiny \textcolor{accentcyan}{\textbf{Scan for Slides/Repo}}}\\
      {\tiny \textcolor{textmuted}{n0ctilucent/talk-ccdc-nov-26}}
    \end{column}
  \end{columns}
  \vfill
  \note{Welcome everyone to tonight's session at the Business School. Direct audience to scan the QR code to pull up the repository and slides on their laptops or phones.}
\end{frame}

%%%%%%%%%%%%%%%%%%%%%%%%%%%%%%%%%%%%%%%%%%%%%%%%%%%%%%%%%%%%%%%%
%% SECTION 1 — Welcome & Agenda
%%%%%%%%%%%%%%%%%%%%%%%%%%%%%%%%%%%%%%%%%%%%%%%%%%%%%%%%%%%%%%%%
\section{Welcome \& Club Agenda}

\begin{frame}{Transamerica Cybersecurity Club}
  \begin{columns}[T]
    \column{0.55\textwidth}
      \textbf{\hl{Club Leadership \& Faculty:}}
      \begin{itemize}
        \item \textbf{Unsu Mwambo} --- President, Transamerica Cybersecurity Club (BSIS '28)
        \item \textbf{Joe Murdock, CISSP, CISM} --- Business School Faculty, BSBA Director, \& NCAE-C PoC
        \item \textbf{Franklin Diaz} --- Technical Session Speaker / Security Architect (\texttt{franklin.diaz@ucdenver.edu})
      \end{itemize}

      \vspace{0.3cm}
      \small\hll{Cyber Defense Center:} \\
      \url{https://www.ucdenver.edu/cyber}

    \column{0.45\textwidth}
      \begin{block}{Club Announcements}
        \begin{itemize}
          \item How to join the club
          \item Spring competition lineup: CCDC, CTFs, \& defenses
          \item Hands-on workshops \& industry guest speakers
        \end{itemize}
      \end{block}
      \vspace{0.25cm}
      \centering
      \footnotesize\textcolor{accentamber}{\textbf{Meeting in JP Morgan Room (1200)}}
  \end{columns}
  \note{Hand off to Unsu for club onboarding, then Joe for CCDC competition timeline.}
\end{frame}

\begin{frame}{Tonight's 90-Minute Flight Plan}
  \begin{description}[style=multiline]
    \item[\hl{5:00 -- 5:15}] \textbf{Welcome \& Club Business} \\ % chktex 8
      Opening remarks, membership onboarding, spring events.
    \item[\hl{5:15 -- 5:35}] \textbf{Act I:\@ The AI Myth \& Next-Token Reality} \\ % chktex 8
      How LLMs actually work: autoregression and subwords.
    \item[\hl{5:35 -- 5:55}] \textbf{The Battle: Hands-On Diabolical Art Contest} \\ % chktex 8
      Students vs.\ models: prompting LLMs to draw ASCII art.
    \item[\hl{5:55 -- 6:10}] \textbf{Act II:\@ Under the Hood of Spatial Failure} \\ % chktex 8
      Tokenization boundaries, 1D vs 2D, and multimodal limits.
    \item[\hl{6:10 -- 6:30}] \textbf{Showdown, Security Implications, \& Extended Q\&A} \\ % chktex 8
      Gallery of horrors, prompt injection risks, 20-min discussion buffer.
  \end{description}
  \note{Pacing note: We built in a dedicated 20-minute buffer at 6:10 to ensure plenty of room for student demo projections and live Q&A.}
\end{frame}

%%%%%%%%%%%%%%%%%%%%%%%%%%%%%%%%%%%%%%%%%%%%%%%%%%%%%%%%%%%%%%%%
%% SECTION 2 — Dispelling the "AI Myth"
%%%%%%%%%%%%%%%%%%%%%%%%%%%%%%%%%%%%%%%%%%%%%%%%%%%%%%%%%%%%%%%%
\section{The ``AI Myth'' vs. Reality}

\begin{frame}{The Illusion of Omniscience}
  \begin{itemize}
    \item \textbf{The Perception:} Modern LLMs write poetry, debug kernel code, pass bar exams, and hold nuanced conversations.
    \item \textbf{The Anthropomorphic Trap:} Because the output is grammatically fluent, we instinctively assume the model possesses \hl{internal reasoning}, a \hl{mental model}, or \hl{visual spatial awareness}.
    \item \textbf{The Reality:} An LLM is an autoregressive token predictor:
      \[
        P(w_t \mid w_1, w_2, \dots, w_{t-1})
      \]
    \item It does not ``plan'' a picture, nor does it possess a 2D drawing canvas.
  \end{itemize}

  \vspace{0.3cm}
  \begin{quote}
    \small\textit{``If all goes to plan, people won't even realize they are learning until it is far too late!''}% chktex 38
  \end{quote}
  \note{Emphasize that linguistic fluency masks architectural limitations. Grammatical competence does not imply spatial awareness.}
\end{frame}

\begin{frame}{What Exactly Is a Token?}
  \begin{columns}[T]
    \column{0.52\textwidth}
      \begin{itemize}
        \item Models do \textbf{not} read individual letters or pixels.
        \item Text is sliced into subwords via \hl{Byte-Pair Encoding (BPE)} or WordPiece.
        \item Words like \texttt{``cyber''} might be 1 token; rare combinations are fragmented.
        \item \textbf{The Whitespace Hazard:} Leading spaces, tabs, and repeating characters map to inconsistent token IDs!
      \end{itemize}

    \column{0.48\textwidth}
      \begin{exampleblock}{Tokenization Example (\texttt{cl100k\_base} / GPT-4)}
        \small
        \texttt{`` wizard''} $\rightarrow$ \texttt{[12842]} \\
        \texttt{``wizard''} $\rightarrow$ \texttt{[38411]} \\
        \texttt{``\ \ *\ \ ''} $\rightarrow$ \texttt{[451, 882]} \\
        \texttt{``\ \ \ *\ ''} $\rightarrow$ \texttt{[1029, 902]}
      \end{exampleblock}
      \vspace{0.2cm}
      \footnotesize\textcolor{accentamber}{\textbf{Takeaway:}} The model cannot count characters when its input vocabulary lumps characters into arbitrary chunks!
  \end{columns}
  \note{Point out how a leading space changes the entire token ID from 38411 to 12842. To a model, ``wizard'' and `` wizard'' are completely different mathematical symbols.}
\end{frame}

%%%%%%%%%%%%%%%%%%%%%%%%%%%%%%%%%%%%%%%%%%%%%%%%%%%%%%%%%%%%%%%%
%% SECTION 3 — The Diabolical Art Battle
%%%%%%%%%%%%%%%%%%%%%%%%%%%%%%%%%%%%%%%%%%%%%%%%%%%%%%%%%%%%%%%%
\section{The Diabolical Art Battle}

\begin{frame}{Rules of Engagement}
  \begin{block}{Participant Challenge}
    Your mission: coax your favorite generative AI into drawing accurate, structured ASCII art to Franklin's \hl{diabolical specifications}.
  \end{block}

  \vspace{0.3cm}
  \begin{itemize}
    \item \textbf{Permitted Weapons:} ChatGPT (GPT-4o/o1), Claude 3.5 Sonnet, Google Gemini 1.5, DeepSeek, or local Ollama models.
    \item \textbf{Strict Rule:} \textbf{Zero human edits.} No backspacing, no manual whitespace fixing, no copy-paste collage.
    \item \textbf{Submission:} Paste raw output into a monospaced editor or share on the room projector.
    \item \textbf{Prize:} Eternal glory for the least deformed creation!
  \end{itemize}
  \note{Remind students: strictly zero human editing allowed. If the model misaligns a brace or drops a column, that's what gets shown on the screen.}
\end{frame}

\begin{frame}{Challenge 1: The Simple Geometric Trap}
  \begin{columns}[T]
    \column{0.5\textwidth}
      \textbf{\hl{Specification:}}
      \begin{itemize}
        \item Construct an exact \textbf{10 $\times$ 10} square box made of asterisks \texttt{*}.
        \item The inside must be empty whitespace.
        \item Place a single capital \texttt{'X'} dead center (row 5, column 5).
        \item Include line character counts.
      \end{itemize}

    \column{0.5\textwidth}
      \textbf{\hlm{What Actually Happens:}}
      \begin{itemize}
        \item Row 1: 10 asterisks
        \item Row 2: 8 asterisks + 3 spaces
        \item Row 3: 11 characters
        \item Center \texttt{X}: Off by 3 columns
      \end{itemize}
      \vspace{0.3cm}
      \centering
      \small\textcolor{accentcyan}{\textbf{A task a 5-line Python script solves in 1 ms defeats a \$100B model.}}
  \end{columns}
  \note{Walk the room as students run Challenge 1. Ask for hands: did anyone's model generate an actual 10 $\times$ 10 square on the first prompt?}
\end{frame}

\begin{frame}[fragile]{Challenge 2: The Infamous ASCII Wizard}
  \begin{columns}[T]
    \column{0.55\textwidth}
      \textbf{\hl{The Diabolical Wizard Prompt:}}
      \begin{enumerate}
        \item A pointed wizard hat with a recognizable brim.
        \item A face with eyes and a long flowing beard.
        \item Robes with two distinct arms.
        \item Right hand holding a magical staff.
        \item Left hand holding an illuminated orb.
      \end{enumerate}

      \vspace{0.25cm}
      \footnotesize\textcolor{accentamber}{\textit{Franklin's Note: ``I was trying to find some of these absolutely terrible wizard drawings I had OpenClaw do\dots{} total eldritch horrors!''}}% chktex 38

    \column{0.45\textwidth}
      \begin{exampleblock}{Typical Model Output}
\begin{lstlisting}
    /\
   /  \
  /====\
 ( . . )  *---Staff
 ~ | | ~ /
  /  |  \
  ===Orb===
\end{lstlisting}
      \end{exampleblock}
      \vspace{0.1cm}
      \footnotesize The hat floats, the staff is detached, and the beard is inside out!
  \end{columns}
  \note{This is the crowd favorite. Model responses will be hilarious; showcase 2--3 examples on the big screen from students.}
\end{frame}

%%%%%%%%%%%%%%%%%%%%%%%%%%%%%%%%%%%%%%%%%%%%%%%%%%%%%%%%%%%%%%%%
%% SECTION 4 — Under the Hood: Why LLMs Fail
%%%%%%%%%%%%%%%%%%%%%%%%%%%%%%%%%%%%%%%%%%%%%%%%%%%%%%%%%%%%%%%%
\section{Why LLMs Fail at ASCII Art}

\begin{frame}{1D Causal Attention vs. 2D Planar Space}
  \begin{itemize}
    \item \textbf{Language is 1-Dimensional:} Words follow words linearly in time:
      \[
        t_1 \longrightarrow t_2 \longrightarrow t_3 \longrightarrow \dots \longrightarrow t_N
      \]
    \item \textbf{Art is 2-Dimensional:} Visual features require spatial alignment across coordinate axes $(x, y)$.
    \item \textbf{No Retraction Buffer:} When an LLM outputs a character on line 4, it cannot move the cursor up to adjust line 2.
    \item \textbf{Cumulative Drift:} A single misplaced whitespace token on column 3 shifts every remaining element on that line and cascades to subsequent lines.
  \end{itemize}
  \note{Deep dive on the absence of a retraction buffer: human artists constantly erase and look back. Autoregressive models must commit irreversibly token by token.}
\end{frame}

\begin{frame}{The Spatial Blind Spot}
  \begin{columns}[T]
    \column{0.5\textwidth}
      \textbf{\hl{Why Spatial Reasoning Fails:}}
      \begin{itemize}
        \item \textbf{Variable Token Width:} A token can represent 1 character (\texttt{``\textbackslash{}''}) or 6 characters (\texttt{``------''}).
        \item \textbf{Attention Distance:} The character directly ``above'' on row $N-1$ is separated by dozens of intermediate tokens in 1D sequence memory.
        \item \textbf{Monospace Blindness:} Training data predominantly mixes variable-width font assumptions.
      \end{itemize}

    \column{0.5\textwidth}
      \begin{block}{The 3eon \& Medium Studies}
        \small
        \begin{itemize}
          \item Even fine-tuning on thousands of ASCII examples yields marginal gain.
          \item Models memorize common templates (e.g.\ basic smiley) but cannot adapt geometry dynamically.
        \end{itemize}
      \end{block}
      \vspace{0.25cm}
      \footnotesize\hll{\url{https://medium.com/data-science/why-llms-suck-at-ascii-art-a9516cb880d5}}
  \end{columns}
  \note{Notice how attention spans degrade over distance: row N-1 is 50--80 tokens away in causal sequence, so maintaining vertical line coordinates is mathematically improbable.}
\end{frame}

\begin{frame}[t]{Can Multi-Modal (Vision) Models Save Us?}
  \begin{columns}[T]
    \column{0.55\textwidth}
      \textbf{\hl{SkyPilot Research:}} \\
      \textit{\small ``Can Multi-Modal LLMs Truly See Images?''}% chktex 38
      \begin{itemize}
        \small
        \item Vision-Language Models (VLMs) patchify images ($14 \times 14$ or $16 \times 16$ grid patches).
        \item They recognize semantic concepts (\textit{``this is a dog''}, \textit{``this is a chart''}).
        \item They struggle with fine-grained spatial coordinates, line intersections, and counting!
        \item Asking a VLM to fix ASCII art generates confident hallucinations of alignment.
      \end{itemize}

    \column{0.45\textwidth}
      \begin{alertblock}{The Vision Paradox}
        \small
        A multi-modal model can describe a picture of a cat, but cannot verify whether column 14 aligns with column 14 on the previous row.
      \end{alertblock}
      \vspace{0.2cm}
      \centering
      \scriptsize\hll{\url{https://skypilot.ai/blog/can-multi-modal-llms-truely-see-images}}
  \end{columns}
  \note{Explain ViT patchification: a 16 $\times$ 16 pixel patch lumps multiple ASCII characters into a single vision token embedding. Fine pixel coordinates are lost.}
\end{frame}

%%%%%%%%%%%%%%%%%%%%%%%%%%%%%%%%%%%%%%%%%%%%%%%%%%%%%%%%%%%%%%%%
%% SECTION 5 — Cybersecurity Implications
%%%%%%%%%%%%%%%%%%%%%%%%%%%%%%%%%%%%%%%%%%%%%%%%%%%%%%%%%%%%%%%%
\section{Cybersecurity Implications}

\begin{frame}{Why Cyber Defenders Should Care}
  \begin{itemize}
    \item \textbf{ArtPrompt / ASCII Smuggling:}
      \begin{itemize}
        \item Attackers exploit safety filter blind spots by rendering prohibited words (e.g., exploit names, passwords, attack commands) as ASCII art.
        \item Safety guardrails evaluate semantic text tokens, failing to recognize the visual payload formed by the ASCII contours!
      \end{itemize}
    \item \textbf{Token Smuggling \& WAF Bypass:}
      \begin{itemize}
        \item Exploiting whitespace and BPE splits to sneak payloads past LLM-based security gateways.
      \end{itemize}
    \item \textbf{Log \& Telemetry Parsing Pitfalls:}
      \begin{itemize}
        \item Never trust an LLM to accurately parse character-aligned logs (hex dumps, packet captures, memory tables) without strict deterministic tooling.
      \end{itemize}
  \end{itemize}
  \note{Tie back to CCDC & defense: never delegate network packet dissection or hex dump boundary checks to an LLM without deterministic validators.}
\end{frame}

\begin{frame}{Defensive Takeaways for Security Engineers}
  \begin{columns}[T]
    \column{0.5\textwidth}
      \begin{block}{What LLMs Excel At}
        \begin{itemize}
          \item Summarizing threat reports
          \item Generating boilerplate scripts
          \item Explaining high-level attack patterns
          \item Natural language query interfaces
        \end{itemize}
      \end{block}

    \column{0.5\textwidth}
      \begin{alertblock}{What Requires Determinism}
        \begin{itemize}
          \item Exact byte / packet parsing
          \item Memory offset calculations
          \item Coordinate / spatial layout auditing
          \item Cryptographic signature checks
        \end{itemize}
      \end{alertblock}
  \end{columns}

  \vspace{0.4cm}
  \centering
  \small\textcolor{accentamber}{\textbf{Principle: Use AI for semantic synthesis; use deterministic code for byte precision.}}
  \note{Core rule of thumb to leave students with: Semantic synthesis = AI;\@ Byte/coordinate precision = Deterministic Code.}
\end{frame}

%%%%%%%%%%%%%%%%%%%%%%%%%%%%%%%%%%%%%%%%%%%%%%%%%%%%%%%%%%%%%%%%
%% SECTION 6 — References & Closing
%%%%%%%%%%%%%%%%%%%%%%%%%%%%%%%%%%%%%%%%%%%%%%%%%%%%%%%%%%%%%%%%
\section{References \& Community}

\begin{frame}{References \& Further Reading}
  \begin{itemize}
    \item \textbf{Experiment 8: LLM-Generated ASCII Art} (3eon): \\
      \footnotesize\url{https://3eon.medium.com/experiment-8-llm-generated-ascii-art-879e96053867}
    \item \textbf{Why LLMs Suck at ASCII Art} (Towards Data Science): \\
      \footnotesize\url{https://medium.com/data-science/why-llms-suck-at-ascii-art-a9516cb880d5}
    \item \textbf{Can Multi-Modal LLMs Truly See Images?} (SkyPilot AI): \\
      \footnotesize\url{https://skypilot.ai/blog/can-multi-modal-llms-truely-see-images}
    \item \textbf{ArtPrompt: ASCII Art-based Jailbreak Attacks against LLMs} (IEEE S\&P / arXiv:2402.11753): \\
      \footnotesize\url{https://arxiv.org/abs/2402.11753}
    \item \textbf{CU Denver Cyber Defense Center}: \\
      \footnotesize\url{https://www.ucdenver.edu/cyber}
  \end{itemize}
  \note{Direct students to the arXiv ArtPrompt paper if they want to study jailbreak mechanics for upcoming CCDC and CTF challenges.}
\end{frame}

%%%%%%%%%%%%%%%%%%%%%%%%%%%%%%%%%%%%%%%%%%%%%%%%%%%%%%%%%%%%%%%%
%% CLOSING SLIDE
%%%%%%%%%%%%%%%%%%%%%%%%%%%%%%%%%%%%%%%%%%%%%%%%%%%%%%%%%%%%%%%%
\begin{frame}[plain]
  \vfill
  \begin{columns}[c]
    \begin{column}{0.72\textwidth}
      {\Huge\bfseries\textcolor{accentcyan}{Thank You!}}\\[0.25cm]
      {\Large\textcolor{white}{Questions, Live Demos, \& Gallery Submissions}}\\[0.5cm]

      \begin{block}{Contact \& Presentation Repository}
        \small
        \textbf{Franklin Diaz} --- \texttt{franklin.diaz@ucdenver.edu} \quad (\texttt{@n0ctilucent}) \\
        \textbf{Transamerica Cybersecurity Club} --- CU Denver Business School \\
        \textbf{Cyber Defense Center (NCAE-C):} \url{https://www.ucdenver.edu/cyber} \\
        \textbf{Slide Repository:} \url{https://github.com/n0ctilucent/talk-ccdc-nov-26}
      \end{block}

      \vspace{0.3cm}
      {\footnotesize \textcolor{textmuted}{Special thanks to Unsu Mwambo, Prof.\ Joe Murdock, and the Cyber Defense Center!}}
    \end{column}
    \begin{column}{0.25\textwidth}
      \centering
      \colorbox{white}{\qrcode[height=2.5cm]{https://github.com/n0ctilucent/talk-ccdc-nov-26}}\\[0.2cm]
      {\tiny \textcolor{accentcyan}{\textbf{Slide Deck \& Code Repo}}}\\
      {\tiny \textcolor{textmuted}{github.com/n0ctilucent/talk-ccdc-nov-26}}
    \end{column}
  \end{columns}
  \vfill
  \note{Open floor for questions and project the top 3 student ASCII wizard submissions.}
\end{frame}

\end{document}
